\documentclass[11pt,a4paper]{article}
\usepackage[utf8]{inputenc}
\usepackage[T1]{fontenc}
\usepackage{times}
\usepackage{geometry}
\geometry{margin=1in}
\usepackage{amsmath,amssymb}
\usepackage{booktabs}
\usepackage{graphicx}
\usepackage{multirow}
\usepackage{xcolor}
\usepackage{hyperref}
\usepackage{natbib}
\usepackage{caption}
\usepackage{subcaption}
\usepackage{tabularx}

\hypersetup{colorlinks=true,linkcolor=blue,citecolor=blue,urlcolor=blue}

\title{Why Lightweight Gating Cannot Fix Multi-Turn Attention Degradation:\\A Mechanistic Analysis}

\author{Anonymous}
\date{}

\begin{document}
\maketitle

% ============================================================
% ABSTRACT
% ============================================================
\begin{abstract}
Multi-turn dialogue performance in large language models degrades as conversations grow longer.
Recent work by \citet{dongre2026} quantified this through the \emph{Goal Attention Ratio} (GAR), showing 27--48\% drops over extended conversations, while residual stream probes reveal that goal information persists (AUC up to 0.98) even as the attention channel closes.
This creates a natural hypothesis: a learnable gate on the residual connection should recover goal attention.
We implement this hypothesis as \textsc{ThreadGate}, a per-layer gating module trained with joint next-token prediction (NTP) and GAR objectives.
ThreadGate successfully optimizes GAR during training ($0.02 \to 0.48$), but this improvement completely fails to transfer to inference (GAR~$= 0.014$, a \textbf{34$\times$} gap on held-out data).
We identify teacher-forcing as the structural cause: attention patterns learned under gold contexts do not generalize to autoregressive generation.
Two coupled mechanisms amplify this shift: (1)~NTP gradients dominate gate parameters with 7--17$\times$ magnitude over GAR gradients, causing gates to learn \emph{reverse} turn conditioning in 27/28 (Qwen) and 30/32 (Llama) layers; (2)~at inference time, gates attenuate residual updates by only 2--7\%, far too weak to shift attention.
Four-way ablation (ThreadGate vs.\ Uniform Gate vs.\ LoRA-only vs.\ Base) on two architectures (Qwen-2.5-7B and Llama-3.1-8B) confirms that gating provides zero incremental contribution over LoRA fine-tuning alone.
Our results demonstrate that auxiliary attention losses optimized under teacher-forcing face a fundamental generalization barrier, with implications for all training-time attention steering methods.
\end{abstract}

% ============================================================
% 1. INTRODUCTION
% ============================================================
\section{Introduction}

Large language models (LLMs) exhibit systematic performance degradation in multi-turn dialogues.
As conversations extend, models progressively lose track of instructions, goals, and context established in early turns---a phenomenon familiar to practitioners but only recently quantified.
\citet{dongre2026} introduced the \emph{Goal Attention Ratio} (GAR), measuring the fraction of attention mass allocated to goal-defining tokens (e.g., system prompt instructions), and demonstrated 27--48\% GAR drops across 12 model families over 50 turns.

Crucially, \citet{dongre2026} also showed that goal information is not \emph{lost} from the model's representations: linear probes on the residual stream achieve AUC scores up to 0.98 for detecting goal-relevant content, even at late turns where attention to goal tokens has collapsed.
This dissociation---information present in residual representations but inaccessible through attention---suggests a natural intervention: a learnable gating mechanism on the residual connection that modulates information flow to restore goal attention.

We implement this hypothesis as \textsc{ThreadGate}, inserting per-layer scalar gates into the residual stream, trained jointly with an auxiliary GAR loss alongside the standard NTP objective.
ThreadGate successfully optimizes training-time GAR from 0.02 to 0.48 over 3 epochs, demonstrating that the auxiliary loss and gating architecture are sound in principle.
However, this improvement \emph{completely fails to transfer to inference}: evaluated on held-out conversations under autoregressive generation, ThreadGate achieves GAR of only 0.014---a 34$\times$ gap from training performance.

Our central finding is that this failure is \emph{structural}, caused by the distribution shift between teacher-forced training and autoregressive inference.
During training, the model conditions on gold (human-written) context when computing attention; at inference, it conditions on its own generated text.
Attention patterns optimized under the former do not transfer to the latter.
We further identify two coupled mechanisms that amplify this shift:
\begin{itemize}
    \item \textbf{Gradient competition}: NTP loss gradients on gate parameters are 7--17$\times$ larger than GAR gradients and directionally opposed in 66--79\% of layers, causing gates to learn reverse turn conditioning---suppressing residual updates at early turns (where GAR is high) and becoming passive at late turns (where intervention is needed).
    \item \textbf{Insufficient intervention strength}: At inference time, learned gates attenuate residual updates by only 2--7\% (gate values 0.93--0.98), far too weak to meaningfully shift attention distributions.
\end{itemize}

Four-way ablation (ThreadGate vs.\ Uniform Gate vs.\ LoRA-only vs.\ Base) shows that gating provides zero incremental GAR improvement over LoRA fine-tuning alone ($<$0.2 percentage points).
Cross-model validation on Llama-3.1-8B (GQA architecture) reproduces the same failure pattern: 32$\times$ training-inference gap with gate direction reversal in 30/32 layers.

Our contributions are:
\begin{enumerate}
    \item \textbf{Quantitative evidence of teacher-forcing distribution shift} destroying auxiliary attention loss generalization (30--34$\times$ training-inference GAR gap), with implications for all training-time attention steering methods.
    \item \textbf{Layer-wise quantification of NTP-GAR gradient competition}, showing NTP magnitude suppression (7--17$\times$) causes gates to learn reverse turn conditioning.
    \item \textbf{Systematic ablation} isolating gating's zero incremental contribution via four-way comparison with dual GAR + perplexity metrics across two architectures.
    \item \textbf{Cross-model validation} (Qwen-2.5-7B transformer + Llama-3.1-8B GQA) confirming the structural nature of the failure.
\end{enumerate}

% ============================================================
% 2. RELATED WORK
% ============================================================
\section{Related Work}

\paragraph{Multi-turn dialogue degradation.}
The observation that LLMs lose coherence over long conversations has been documented across instruction-following benchmarks \citep{zheng2024chatbot,li2024longcontext}.
\citet{dongre2026} provided the first systematic quantification through GAR, demonstrating the phenomenon across 12 model families and linking it to attention distribution shifts.
Their residual stream probes further showed that goal information persists in hidden states even when attention fails to surface it.

\paragraph{Residual gating and adaptive computation.}
Gating mechanisms on residual connections have been explored for adaptive depth \citep{raposo2024mixture,fan2020reducing} and layer selection \citep{elbayad2020depth}.
Our work applies gating not for computational efficiency but for \emph{steering} information flow, training gates with an auxiliary attention objective.

\paragraph{Attention steering and intervention.}
Training-time attention regularization has been used for various objectives \citep{ghader2017attention,strubell2018linguistically}.
Inference-time attention intervention methods \citep{li2023inference,todd2024function} modify attention patterns without training, bypassing the distribution shift we identify.
Our findings suggest that the latter class of approaches may be more promising for multi-turn goal retention.

\paragraph{Exposure bias.}
The discrepancy between teacher-forced training and autoregressive inference is well-studied for generation quality \citep{ranzato2016sequence,bengio2015scheduled}.
Our work extends this analysis to auxiliary attention losses, providing the first quantitative evidence that exposure bias can create order-of-magnitude gaps in auxiliary objective transfer.

% ============================================================
% 3. BACKGROUND
% ============================================================
\section{Background}

\subsection{Goal Attention Ratio (GAR)}

Following \citet{dongre2026}, we define the Goal Attention Ratio as the fraction of attention mass allocated to goal-defining tokens.
We employ two GAR variants:

\textbf{GAR-Dongre}: Goal tokens are system prompt task tokens only; response tokens are all tokens in assistant turns.
This definition enables direct comparison with \citet{dongre2026}.

\textbf{GAR-Ours}: Goal tokens include both the system prompt and the first user query; response tokens are the last token of each assistant turn (as a summary position).
This provides a richer training signal.

Both definitions yield consistent results: inference GAR is approximately 0.014--0.016 across definitions, forming a 30--34$\times$ gap with training GAR (0.42--0.48).

\subsection{Residual Stream Persistence}

\citet{dongre2026} demonstrated that linear probes trained on residual stream activations achieve AUC up to 0.98 for detecting goal-relevant information at late conversation turns.
This means the model retains goal information in its hidden states even when attention to goal tokens has collapsed---the information exists but the attention mechanism fails to access it.
This dissociation motivates our gating intervention: if goal information is present in the residual stream, modulating the residual connection should be able to surface it through the attention mechanism.

% ============================================================
% 4. METHOD
% ============================================================
\section{Method}

\subsection{ThreadGate Architecture}

We insert a learnable scalar gate $g_l \in (0, 1)$ at each transformer layer $l$.
The gated residual update is:
\begin{equation}
    \mathbf{h}_l = \mathbf{h}_{l-1} + g_l \cdot f_l(\mathbf{h}_{l-1})
\end{equation}
where $f_l$ denotes the combined self-attention and feed-forward computation at layer $l$, and $g_l = \sigma(\alpha_l)$ with $\alpha_l$ a learnable parameter initialized to $\sigma^{-1}(0.95) \approx 2.94$.

Gates are initialized near 1.0 (identity) to preserve pretrained model behavior at the start of training.
Low-Rank Adaptation (LoRA) modules are applied to attention projections ($W_Q, W_K, W_V, W_O$) to provide the model with additional capacity for adapting attention patterns.

\subsection{Training Objective}

The training loss combines next-token prediction with a GAR-based auxiliary objective:
\begin{equation}
    \mathcal{L} = \mathcal{L}_{\text{NTP}} + \lambda_{\text{GAR}} \cdot \mathcal{L}_{\text{GAR}} + \lambda_{\text{div}} \cdot \mathcal{L}_{\text{div}}
\end{equation}

\textbf{GAR loss ($\mathcal{L}_{\text{GAR}}$).}
We use a mean squared error (MSE) target that encourages each layer's GAR to match a target value derived from the conversation structure:
\begin{equation}
    \mathcal{L}_{\text{GAR}} = \frac{1}{L} \sum_{l=1}^{L} \left( \text{GAR}_l - \text{GAR}_{\text{target}} \right)^2
\end{equation}
The MSE formulation is bidirectional---it can both increase and decrease GAR---providing more stable gradients than a pure maximization objective.

\textbf{Gate diversity loss ($\mathcal{L}_{\text{div}}$).}
To prevent gate collapse (all gates converging to the same value), we add a diversity regularizer:
\begin{equation}
    \mathcal{L}_{\text{div}} = -\text{Var}(\{g_l\}_{l=1}^{L})
\end{equation}

\subsection{Loss Design Evolution (v1--v7)}

We explored seven training configurations spanning four core loss families before arriving at the MSE + gate diversity formulation (v7):
\begin{itemize}
    \item \textbf{v1--v2} (cumulative): Accumulated GAR over turns. Failed to converge due to gradient instability.
    \item \textbf{v3--v5} (maximize): Direct GAR maximization. Unstable training with mode collapse.
    \item \textbf{v6} (smooth saturation): Saturating activation on GAR. Partial convergence but poor GAR recovery.
    \item \textbf{v7} (MSE target): Bidirectional MSE loss + gate diversity. Successfully converges; adopted for all experiments.
\end{itemize}

\subsection{Ablation Variants}

To isolate the contribution of gating, we train three additional configurations:
\begin{itemize}
    \item \textbf{Uniform Gate (UG)}: All gates fixed to the same learned scalar (no per-layer or per-turn variation). Isolates per-layer turn conditioning from overall gating.
    \item \textbf{LoRA-only (LO)}: LoRA modules only, gates fixed at 1.0, NTP loss only. Isolates the contribution of LoRA fine-tuning from gating.
    \item \textbf{Base}: Unmodified pretrained model. Establishes the baseline.
\end{itemize}

\subsection{Training Setup}

All models are trained for 3 epochs on multi-turn conversations from UltraChat.
We use the same hyperparameters across architectures: LoRA rank 16, learning rate $2 \times 10^{-4}$, batch size 4 with gradient accumulation.
Qwen-2.5-7B (28 layers, full MHA) and Llama-3.1-8B-Instruct (32 layers, GQA) serve as our two architectures.

% ============================================================
% 5. EXPERIMENTS
% ============================================================
\section{Experiments}

\subsection{Baseline GAR Decay Across Architectures}

We first verify that GAR decay is a cross-architecture phenomenon by measuring per-turn GAR on 5 standard conversations (10 turns each) for both base models.

\begin{table}[h]
\centering
\caption{Cross-model GAR decay in base (unmodified) models. Both architectures exhibit $>$70\% relative GAR drop from turn~1 to mid-turns, confirming the structural nature of attention degradation.}
\label{tab:gar-decay}
\begin{tabular}{lcc}
\toprule
\textbf{Metric} & \textbf{Qwen-2.5-7B} & \textbf{Llama-3.1-8B} \\
\midrule
Turn 1 GAR       & 0.0418 & 0.0355 \\
Mid-turns GAR    & 0.0122 & 0.0098 \\
Relative Drop    & 70.8\% & 72.4\% \\
Decay Shape      & Monotone $\downarrow$ & Monotone $\downarrow$ \\
\bottomrule
\end{tabular}
\end{table}

Both Qwen-2.5-7B (full MHA, 28 layers) and Llama-3.1-8B (GQA, 32 layers) show monotonically decreasing GAR with $>$70\% relative drop (Table~\ref{tab:gar-decay}).
The slightly lower absolute GAR in Llama is consistent with GQA's more uniform attention distribution due to KV sharing, but the decay rate is comparable.
This confirms that GAR decay is architecture-independent, driven by softmax normalization over growing context.

\subsection{Training Dynamics}

ThreadGate training on both architectures exhibits a characteristic three-phase pattern:

\begin{itemize}
    \item \textbf{Phase A} (plateau): GAR remains low ($\sim$0.015) as NTP dominates gradients (Qwen: gs0--90; Llama: gs0--90).
    \item \textbf{Phase B} (rapid rise): GAR rises sharply as the MSE loss finds gradient pathways in low-conflict layers (Qwen: gs90--144; Llama: gs90--170).
    \item \textbf{Phase C} (convergence): GAR stabilizes at 0.42--0.48 as the GAR training signal is exhausted.
\end{itemize}

Qwen converges at global step 144 (training GAR = 0.48); Llama converges later at step 170 (training GAR = 0.42), consistent with its deeper architecture requiring more steps for gate coordination across 32 layers.
Final gate statistics are similar: Qwen gate mean 0.93 ($\sigma$=0.05), Llama gate mean 0.94 ($\sigma$=0.04), both with gates near identity.

\subsection{Four-Way Ablation: Qwen-2.5-7B}

We evaluate all four model variants on 50 held-out conversations (UltraChat idx$\geq$8000, never seen during training) using both GAR and perplexity (PPL) as metrics.

\begin{table}[h]
\centering
\caption{Held-out four-way evaluation on Qwen-2.5-7B (50 conversations, idx$\geq$8000). Gating provides negligible incremental GAR improvement ($<$0.2pp) over LoRA-only, and does not affect perplexity.}
\label{tab:qwen-heldout}
\begin{tabular}{lcccc}
\toprule
\textbf{Model} & \textbf{GAR-Dongre} & \textbf{GAR-Ours} & \textbf{PPL} & \textbf{NTP Loss} \\
\midrule
Base        & 0.0025 & 0.009 & 4.45 & 1.44 \\
ThreadGate  & 0.014  & 0.015 & 2.71 & 0.96 \\
Uniform Gate & 0.012 & 0.016 & 2.69 & 0.95 \\
LoRA-only   & 0.013  & 0.014 & 2.70 & 0.96 \\
\bottomrule
\end{tabular}
\end{table}

Table~\ref{tab:qwen-heldout} reveals that all fine-tuned variants (TG, UG, LO) achieve similar GAR and PPL, with ThreadGate providing only +0.1pp GAR over LoRA-only.
The substantial improvement over Base (GAR: 0.0025$\to$0.013; PPL: 4.45$\to$2.70) is entirely attributable to LoRA fine-tuning on NTP, not to the gating mechanism.

\subsection{Three-Way Evaluation: Llama-3.1-8B}

We replicate the ablation on Llama-3.1-8B with Base, ThreadGate, and LoRA-only on held-out data.

\begin{table}[h]
\centering
\caption{Llama-3.1-8B three-way evaluation on held-out data. The pattern mirrors Qwen: gating provides zero incremental benefit.}
\label{tab:llama-heldout}
\begin{tabular}{lccc}
\toprule
\textbf{Model} & \textbf{GAR-Dongre} & \textbf{PPL} & \textbf{NTP Loss} \\
\midrule
Base        & 0.0022 & 4.52 & 1.46 \\
ThreadGate  & 0.013  & 2.78 & 0.975 \\
LoRA-only   & 0.012  & 2.76 & 0.970 \\
\bottomrule
\end{tabular}
\end{table}

Llama results (Table~\ref{tab:llama-heldout}) confirm the Qwen findings: ThreadGate--LoRA GAR difference is +0.1pp (not significant), and LoRA-only achieves marginally better PPL (2.76 vs.\ 2.78).
Gating is redundant across both architectures.

\subsection{Cross-Model Gating Incremental Benefit}

Table~\ref{tab:incremental} consolidates the gating incremental analysis across all evaluation settings.

\begin{table}[h]
\centering
\caption{Cross-model gating incremental benefit (ThreadGate minus LoRA-only). All differences are within noise across three independent evaluation settings.}
\label{tab:incremental}
\begin{tabular}{lccc}
\toprule
\textbf{Comparison} & \textbf{Qwen (train-overlap)} & \textbf{Qwen (held-out)} & \textbf{Llama (held-out)} \\
\midrule
TG $-$ LO GAR (pp) & +0.20 & +0.10 & +0.10 \\
TG $-$ LO PPL       & 0.00  & +0.01 & +0.02 \\
Significant?         & No    & No    & No \\
\bottomrule
\end{tabular}
\end{table}

\subsection{The Training-Inference Gap}

The central quantitative finding is the dramatic gap between training-time and inference-time GAR performance.

\begin{table}[h]
\centering
\caption{Training-inference gap across architectures. Both models show $>$30$\times$ GAR collapse from training (teacher-forced) to inference (autoregressive), with gate direction reversal in nearly all layers.}
\label{tab:gap}
\begin{tabular}{lcc}
\toprule
\textbf{Metric} & \textbf{Qwen-2.5-7B} & \textbf{Llama-3.1-8B} \\
\midrule
Training GAR (converged)   & 0.48   & 0.42 \\
Inference GAR (Dongre, held-out) & 0.014 & 0.013 \\
Gap                        & 34$\times$ & 32$\times$ \\
Gate wrong direction       & 27/28 layers & 30/32 layers \\
Convergence step           & gs144  & gs170 \\
\bottomrule
\end{tabular}
\end{table}

Table~\ref{tab:gap} shows that despite successful training convergence, inference GAR drops by 32--34$\times$.
This gap is not an implementation artifact (verified by line-by-line code audit and controlled experiments) but a direct consequence of the distribution shift between teacher-forced and autoregressive contexts.

\subsection{Held-out vs.\ Training-Overlap Consistency}

To rule out the possibility that the 30$\times$ gap is an overfitting artifact, we compare results on training-overlap and held-out evaluation data.

\begin{table}[h]
\centering
\caption{Consistency between training-overlap and held-out evaluation (Qwen-2.5-7B). The held-out gap is slightly \emph{larger} (34$\times$ vs.\ 30$\times$), ruling out data memorization.}
\label{tab:consistency}
\begin{tabular}{lccc}
\toprule
\textbf{Metric} & \textbf{Train-Overlap} & \textbf{Held-out} & \textbf{Consistent?} \\
\midrule
TG Dongre mid GAR     & 0.0157 & 0.014 & Yes \\
Training-Inference Gap & 30$\times$   & 34$\times$  & Yes \\
TG--LO GAR diff (pp)  & 0.20   & 0.10  & Yes ($<$0.6pp) \\
TG PPL                 & 2.63   & 2.71  & Yes \\
\bottomrule
\end{tabular}
\end{table}

The held-out gap (34$\times$) is slightly larger than the training-overlap gap (30$\times$), consistent with greater distribution shift on unseen conversation styles.
This rules out data memorization as an explanation.

% ============================================================
% 6. ANALYSIS
% ============================================================
\section{Analysis}

\subsection{NTP-GAR Gradient Competition}

To understand \emph{why} gates fail to learn useful turn conditioning, we analyze the gradient dynamics on gate parameters.
At each epoch checkpoint, we compute the NTP and GAR loss gradients with respect to each gate parameter $\alpha_l$ and measure: (1)~cosine similarity (directional alignment) and (2)~magnitude ratio (relative strength).

\begin{table}[h]
\centering
\caption{NTP-GAR gradient conflict evolution during training (Qwen-2.5-7B). NTP magnitude suppression increases monotonically while directional conflict persists in the majority of layers.}
\label{tab:grad-qwen}
\begin{tabular}{lcccc}
\toprule
\textbf{Checkpoint} & \textbf{\% Layers cos$<$0} & \textbf{Mean cos\_sim} & \textbf{Median mag\_ratio} & \textbf{L27 mag\_ratio} \\
\midrule
Epoch 0 (gs1000) & 79\% & $-$0.089 & 7.5$\times$  & 135$\times$ \\
Epoch 1 (gs2000) & 71\% & $-$0.062 & 11.2$\times$ & 580$\times$ \\
Epoch 2 (gs3000) & 67\% & $-$0.042 & 14.7$\times$ & 1139$\times$ \\
\bottomrule
\end{tabular}
\end{table}

Three findings emerge from Table~\ref{tab:grad-qwen}:
\begin{enumerate}
    \item \textbf{Persistent directional conflict}: 67--79\% of layers show opposing NTP and GAR gradient directions throughout training. The gradual decrease reflects GAR signal marginalization, not reconciliation.
    \item \textbf{Monotonic magnitude suppression}: NTP gradient magnitude grows from 7.5$\times$ to 14.7$\times$ over GAR, meaning gate updates are dominated by NTP.
    \item \textbf{Extreme deep-layer suppression}: Layer 27 (deepest) reaches 1139$\times$ magnitude ratio, explaining why the deepest gates converge to near-identity ($g_{27} \approx 1.0$).
\end{enumerate}

This gradient competition directly explains the three-phase training dynamics:
Phase~A corresponds to high conflict + moderate suppression (gates trapped in NTP local optimum);
Phase~B occurs when GAR loss finds breakthrough pathways in low-conflict layers;
Phase~C reflects GAR training signal exhaustion as magnitude ratios grow too large.

\subsection{Cross-Model Gradient Conflict}

We replicate the gradient analysis on Llama-3.1-8B to verify the pattern is architecture-independent.

\begin{table}[h]
\centering
\caption{Cross-model gradient conflict comparison at final checkpoint. Both architectures show nearly identical conflict patterns despite different attention mechanisms (full MHA vs.\ GQA).}
\label{tab:grad-cross}
\begin{tabular}{lcc}
\toprule
\textbf{Metric} & \textbf{Qwen-2.5-7B (28L)} & \textbf{Llama-3.1-8B (32L)} \\
\midrule
\% Layers cos$<$0    & 67\%   & 66\% \\
Mean cos\_sim        & $-$0.042 & $-$0.048 \\
Median mag\_ratio    & 14.7$\times$ & 16.3$\times$ \\
Deepest layer ratio  & 1139$\times$ (L27) & 980$\times$ (L31) \\
Trend: cos\_sim      & $\uparrow$ toward 0 & $\uparrow$ toward 0 \\
Trend: mag\_ratio    & $\uparrow$ monotone  & $\uparrow$ monotone \\
\bottomrule
\end{tabular}
\end{table}

The cross-model consistency (Table~\ref{tab:grad-cross}) is striking: directional conflict (66--67\%), magnitude suppression trends, and deep-layer extreme ratios are all reproduced.
The slightly stronger initial conflict in Llama ($-$0.105 at epoch~0 vs.\ $-$0.089 for Qwen) is consistent with GQA's shared KV heads creating more entangled gradient pathways.
This confirms that NTP-GAR gradient competition is an inherent property of teacher-forcing multi-task training, not an architecture-specific artifact.

\subsection{Attention Visualization}

We extract per-layer attention patterns for all four Qwen model variants across 3 conversations, computing per-layer GAR (attention mass on goal tokens) at turn~10.

\begin{table}[h]
\centering
\caption{Mean layer GAR at turn~10 under autoregressive inference (Qwen-2.5-7B). All four models are indistinguishable at inference time, despite ThreadGate achieving high GAR during training.}
\label{tab:attn}
\begin{tabular}{lcccc}
\toprule
\textbf{Model} & \textbf{Shallow (L0--L9)} & \textbf{Mid (L10--L19)} & \textbf{Deep (L20--L27)} & \textbf{All Layers} \\
\midrule
Base        & 0.012 & 0.007 & 0.005 & 0.008 \\
ThreadGate  & 0.016 & 0.011 & 0.008 & 0.012 \\
Uniform Gate & 0.014 & 0.010 & 0.007 & 0.011 \\
LoRA-only   & 0.015 & 0.011 & 0.008 & 0.012 \\
\bottomrule
\end{tabular}
\end{table}

Table~\ref{tab:attn} confirms that all fine-tuned models are indistinguishable in their attention patterns at inference.
The contrast with training-time attention is dramatic: under teacher-forcing, ThreadGate's attention patterns show strong goal-token bias, but this steering completely vanishes under autoregressive generation.
Deep layers show faster GAR decay (0.005--0.008 at L20--L27 vs.\ 0.012--0.016 at L0--L9), consistent with deeper layers attending more to local context.

\subsection{Downstream Task Correlation}

To validate that GAR is not merely an abstract metric, we design a simple instruction retention test: the system prompt contains an explicit instruction (e.g., ``always begin responses with [OK]''), and we measure compliance rate over 20 turns.

\begin{table}[h]
\centering
\caption{Instruction retention rate by turn. All models show progressive instruction forgetting, with ThreadGate providing no significant improvement over LoRA-only or Base.}
\label{tab:retention}
\begin{tabular}{lccc}
\toprule
\textbf{Turn} & \textbf{Base} & \textbf{ThreadGate} & \textbf{LoRA-only} \\
\midrule
1   & 95\% & 96\% & 96\% \\
5   & 72\% & 74\% & 73\% \\
10  & 48\% & 50\% & 49\% \\
20  & 25\% & 27\% & 26\% \\
\midrule
Overall & 55\% & 57\% & 56\% \\
\bottomrule
\end{tabular}
\end{table}

Instruction retention drops from 95\% at turn~1 to 25\% at turn~20 for the base model (Table~\ref{tab:retention}), and ThreadGate provides only +1--2pp improvement (not statistically significant).
The Pearson correlation between per-turn-bucket GAR and retention rate is $r = 0.82$, confirming that GAR captures a real functional degradation: as attention to goal tokens decreases, the model progressively forgets its instructions.

\begin{table}[h]
\centering
\caption{GAR-retention correlation by turn bucket (Base model). The strong correlation ($r=0.82$) validates GAR as a proxy for functional instruction adherence.}
\label{tab:gar-retention}
\begin{tabular}{lccc}
\toprule
\textbf{Turn Bucket} & \textbf{Mean GAR} & \textbf{Mean Retention} \\
\midrule
Turn 1--5   & 0.030 & 0.84 \\
Turn 6--10  & 0.009 & 0.55 \\
Turn 11--20 & 0.005 & 0.30 \\
\midrule
\multicolumn{2}{l}{Overall Pearson $r$} & 0.82 \\
\bottomrule
\end{tabular}
\end{table}

% ============================================================
% 7. DISCUSSION
% ============================================================
\section{Discussion}

\subsection{Why Teacher-Forcing Breaks Auxiliary Attention Losses}

Our results reveal a fundamental limitation of training-time attention steering under teacher-forcing.
The 30--34$\times$ training-inference GAR gap is not specific to our gating architecture or loss design---it arises from the distributional mismatch between gold and generated contexts.
During training, the model computes attention over human-written text, and the auxiliary loss successfully steers attention toward goal tokens in this regime.
At inference, the model conditions on its own generations, which differ in style, length, and content distribution from the training context.
Attention patterns are sensitive to these distributional differences, and patterns learned under one distribution do not transfer to the other.

This finding has broad implications.
\citet{dongre2026} independently reported failure of goal re-injection (their Appendix~H), where injecting goal representations into later turns did not improve GAR at inference.
Combined with our gating failure, two independent intervention strategies have now encountered the same barrier.
We argue that \emph{any} auxiliary attention loss optimized under teacher-forcing will face this fundamental generalization problem.

\subsection{The Gradient Competition Trap}

The gradient competition we quantify (Section~6.1) reveals a second, coupled failure mode.
Even if the distribution shift were somehow mitigated, NTP's 7--17$\times$ gradient magnitude advantage means gate parameters are effectively controlled by NTP rather than GAR.
The resulting reverse turn conditioning (gates suppressing at early turns, passive at late turns) is the \emph{opposite} of the desired behavior.
This is a classical multi-task learning failure: when task losses differ by an order of magnitude, the weaker task's signal is marginalized.
Gradient balancing techniques \citep{chen2018gradnorm} could potentially mitigate this, but would not address the underlying distribution shift.

\subsection{Constructive Implications}

Our negative results point toward several promising research directions:

\textbf{Inference-time attention intervention.}
Methods that modify attention patterns at inference time \citep{li2023inference} bypass the training-inference distribution shift entirely, as they operate directly on the generation-time context.

\textbf{Decoupled goal-processing pathways.}
Rather than steering the main attention mechanism, a separate processing pathway could maintain goal representations independently of the growing conversation context.

\textbf{GAR-aware decoding.}
Incorporating GAR monitoring into the decoding process (e.g., re-attending to goal tokens when GAR drops below a threshold) could provide dynamic correction without requiring training-time optimization.

\textbf{Exposure bias mitigation for auxiliary losses.}
Techniques such as scheduled sampling \citep{bengio2015scheduled} or sequence-level training \citep{ranzato2016sequence} could be adapted to reduce the distributional gap for auxiliary attention objectives.

% ============================================================
% 8. CONCLUSION
% ============================================================
\section{Conclusion}

We investigated whether learnable residual gating can recover goal attention in multi-turn LLM dialogues.
Despite successful training-time GAR optimization (0.02$\to$0.48), ThreadGate's improvement completely fails to transfer to inference, with a 30--34$\times$ GAR gap across two architectures (Qwen-2.5-7B and Llama-3.1-8B) and evaluation settings (training-overlap and held-out).
Four-way ablation confirms that gating provides zero incremental contribution over LoRA fine-tuning alone.

Our mechanistic analysis identifies teacher-forcing distribution shift as the structural cause, amplified by NTP-GAR gradient competition (7--17$\times$ magnitude suppression, 66--79\% directional conflict) that causes gates to learn reverse turn conditioning.
These findings generalize beyond our specific architecture: any auxiliary attention loss optimized under teacher-forcing faces the same distribution shift barrier at inference time.

We hope this analysis serves as both a cautionary finding for the attention steering community and a catalyst for exploring inference-time and exposure-bias-aware approaches to the multi-turn goal retention problem.

\section*{Limitations}

Our study has several limitations.
First, we evaluate on two model architectures (Qwen-2.5-7B and Llama-3.1-8B) at the 7--8B scale; behavior at larger scales remains an open question, though the structural nature of the distribution shift argument suggests it will persist.
Second, our downstream evaluation is limited to a simple instruction retention task; more comprehensive behavioral benchmarks would strengthen the practical relevance of the GAR metric.
Third, we do not experiment with exposure bias mitigation techniques (e.g., scheduled sampling applied to the auxiliary loss), which could partially bridge the training-inference gap.
Finally, our training data (UltraChat) represents a single conversation domain; verification across diverse dialogue datasets would further strengthen generalizability claims.

% ============================================================
% REFERENCES
% ============================================================
\bibliographystyle{plainnat}

\begin{thebibliography}{20}

\bibitem[Bengio et~al.(2015)]{bengio2015scheduled}
Samy Bengio, Oriol Vinyals, Navdeep Jaitly, and Noam Shazeer.
\newblock Scheduled sampling for sequence prediction with recurrent neural networks.
\newblock In \emph{Advances in Neural Information Processing Systems}, 2015.

\bibitem[Chen et~al.(2018)]{chen2018gradnorm}
Zhao Chen, Vijay Badrinarayanan, Chen-Yu Lee, and Andrew Rabinovich.
\newblock GradNorm: Gradient normalization for adaptive loss balancing in deep multitask networks.
\newblock In \emph{Proceedings of ICML}, 2018.

\bibitem[Dongre et~al.(2026)]{dongre2026}
Aditya Dongre, Abhay Kumar, and colleagues.
\newblock Goal attention ratio: Quantifying multi-turn attention degradation in large language models.
\newblock In \emph{Proceedings of NAACL}, 2026.

\bibitem[Elbayad et~al.(2020)]{elbayad2020depth}
Maha Elbayad, Jiatao Gu, Edouard Grave, and Michael Auli.
\newblock Depth-adaptive transformer.
\newblock In \emph{Proceedings of ICLR}, 2020.

\bibitem[Fan et~al.(2020)]{fan2020reducing}
Angela Fan, Edouard Grave, and Armand Joulin.
\newblock Reducing transformer depth on demand with structured dropout.
\newblock In \emph{Proceedings of ICLR}, 2020.

\bibitem[Ghader and Monz(2017)]{ghader2017attention}
Hamidreza Ghader and Christof Monz.
\newblock What does attention in neural machine translation pay attention to?
\newblock In \emph{Proceedings of the Eighth IJCNLP}, 2017.

\bibitem[Li et~al.(2023)]{li2023inference}
Kenneth Li, Oam Patel, Fernanda Vi{\'e}gas, Hanspeter Pfister, and Martin Wattenberg.
\newblock Inference-time intervention: Eliciting truthful answers from a language model.
\newblock In \emph{Advances in Neural Information Processing Systems}, 2023.

\bibitem[Li et~al.(2024)]{li2024longcontext}
Dacheng Li and colleagues.
\newblock How long can LLMs maintain coherent context in multi-turn dialogue?
\newblock \emph{arXiv preprint}, 2024.

\bibitem[Ranzato et~al.(2016)]{ranzato2016sequence}
Marc'Aurelio Ranzato, Sumit Chopra, Michael Auli, and Wojciech Zaremba.
\newblock Sequence level training with recurrent neural networks.
\newblock In \emph{Proceedings of ICLR}, 2016.

\bibitem[Raposo et~al.(2024)]{raposo2024mixture}
David Raposo, Sam Ritter, Blake Richards, Timothy Lillicrap, Peter~Conway Humphreys, and Adam Santoro.
\newblock Mixture-of-depths: Dynamically allocating compute in transformer-based language models.
\newblock In \emph{Proceedings of ICML}, 2024.

\bibitem[Strubell et~al.(2018)]{strubell2018linguistically}
Emma Strubell, Patrick Verga, Daniel Andor, David Weiss, and Andrew McCallum.
\newblock Linguistically-informed self-attention for semantic role labeling.
\newblock In \emph{Proceedings of EMNLP}, 2018.

\bibitem[Todd et~al.(2024)]{todd2024function}
Eric Todd, Millicent~Li Li, Arnab~Sen Sharma, Aaron Mueller, Byron~C. Wallace, and David Bau.
\newblock Function vectors in large language models.
\newblock In \emph{Proceedings of ICLR}, 2024.

\bibitem[Zheng et~al.(2024)]{zheng2024chatbot}
Lianmin Zheng, Wei-Lin Chiang, Ying Sheng, Siyuan Zhuang, Zhanghao Wu, Yonghao Zhuang, Zi~Lin, Zhuohan Li, Dacheng Li, Eric~P. Xing, Hao Zhang, Joseph~E. Gonzalez, and Ion Stoica.
\newblock Judging LLM-as-a-Judge with MT-Bench and Chatbot Arena.
\newblock In \emph{Advances in Neural Information Processing Systems}, 2024.

\end{thebibliography}

\end{document}
